% ============================================================================
%  Project Report — Template Format
%  Deep Learning-Based Time-of-Day Estimation from Sky Images
%  and EXIF Calendar Dates
% ============================================================================
\documentclass[9pt, a4paper, twocolumn]{article}

% ── Packages ────────────────────────────────────────────────────────────────
\usepackage[utf8]{inputenc}
\usepackage[T1]{fontenc}
\usepackage{mathptmx}            % Serif font similar to Cambria
\usepackage[left=2cm, right=2cm, top=2.5cm, bottom=2.5cm]{geometry}
\usepackage{amsmath, amssymb}
\usepackage{graphicx}
\usepackage{booktabs}
\usepackage{array}
\usepackage{float}
\usepackage{hyperref}
\usepackage[english]{babel}
\usepackage{setspace}
\usepackage{titlesec}
\usepackage{enumitem}
\usepackage{caption}
\usepackage{fancyhdr}
\usepackage{multicol}
\usepackage{tabularx}

% ── Spacing ─────────────────────────────────────────────────────────────────
\singlespacing
\setlength{\parindent}{0pt}
\setlength{\parskip}{5pt}
\setlength{\columnsep}{0.6cm}

% ── Section formatting (bold, 9pt) ─────────────────────────────────────────
\titleformat{\section}
  {\normalsize\bfseries}{\thesection.}{0.5em}{}
\titleformat{\subsection}
  {\normalsize\bfseries}{\thesubsection.}{0.5em}{}
\titlespacing*{\section}{0pt}{10pt}{5pt}
\titlespacing*{\subsection}{0pt}{8pt}{4pt}

% ── Caption formatting ─────────────────────────────────────────────────────
\captionsetup{font=small, labelfont=bf, skip=5pt}

% ── Hyperref setup ─────────────────────────────────────────────────────────
\hypersetup{
    colorlinks=true,
    linkcolor=black,
    citecolor=black,
    urlcolor=blue
}

% ============================================================================
\begin{document}

% ── Title (14pt, left-aligned) ─────────────────────────────────────────────
\twocolumn[
\begin{@twocolumnfalse}
\begin{flushleft}
{\fontsize{14}{17}\selectfont\bfseries
Deep Learning-Based Time-of-Day Estimation from Sky Images and EXIF Calendar Dates}
\vspace{8pt}

{\normalsize
Alk\i m G\"onen\c{c} Efe$^{1}$, Damla Parlaky\i ld\i z$^{1}$, Sarp S\"unb\"ul$^{1}$}
\vspace{4pt}

{\small\itshape
$^{1}$Department of Computer Engineering, Gebze Technical University, Kocaeli, T\"urkiye}
\vspace{4pt}

\rule{\textwidth}{0.4pt}

\vspace{4pt}
{\small\bfseries Abstract}
\vspace{2pt}

{\small
Automated image organization systems in modern digital galleries frequently rely on simplistic brightness-based heuristics to classify outdoor photographs by time of day, leading to systematic misclassifications under artificial lighting and overcast conditions. This paper presents a hybrid deep learning framework that estimates the exact time of capture from a single sky image by fusing visual features extracted by a ConvNeXt-Tiny convolutional backbone with cyclically encoded EXIF calendar metadata. The proposed architecture integrates 77-dimensional handcrafted photometric descriptors---including color-space statistics, sky-region luminance gradients, and edge density measures---with learned deep representations through a metadata fusion multilayer perceptron. A cyclic sine--cosine target encoding is employed to model the circular nature of the 24-hour clock, ensuring continuity at the midnight boundary. The training pipeline incorporates Mixup data augmentation, label noise injection, cosine annealing scheduling, and weighted sampling to address class imbalance across temporal bins. Additionally, a parallel MATLAB-based pipeline utilizing SqueezeNet with Random Forest and Support Vector Regression baselines is developed for comparative analysis. Five-fold cross-validation is used for model evaluation.
}
\vspace{4pt}

{\small\bfseries Keywords:} {\small time-of-day estimation, convolutional neural networks, ConvNeXt, EXIF metadata, cyclic regression, sky image analysis, feature fusion}

\vspace{6pt}
\rule{\textwidth}{0.4pt}
\vspace{12pt}
\end{flushleft}
\end{@twocolumnfalse}
]

% ========================================================================
%  1. INTRODUCTION
% ========================================================================
\section{Introduction}

The estimation of time of day from a single outdoor photograph is a challenging problem at the intersection of computer vision, digital image processing, and environmental sensing. Modern mobile operating systems and cloud storage services attempt to semantically organize landscape photographs, yet their algorithms frequently produce unreliable temporal classifications due to an over-reliance on simplistic brightness metrics [1].

Two principal failure modes motivate this research. First, nighttime images containing strong artificial light sources---such as street lamps, stadium floodlights, or illuminated building facades---are routinely misclassified as ``Daylight'' or ``Sunny'' because high localized luminance inflates the average pixel intensity. Second, midday photographs captured under heavy cloud cover or in deeply shaded areas are often mislabeled as ``Night'' or ``Twilight'' due to globally suppressed light levels that mislead intensity-based classifiers.

These systematic errors underscore a fundamental limitation: average brightness alone is an insufficient proxy for solar position. A robust estimator must instead analyze complex visual cues---shadow geometry, atmospheric color temperature shifts from blue to warm orange, sky gradient profiles, and textural patterns in cloud formations---to replicate the human cognitive ability to distinguish natural from artificial illumination [2].

A foundational study by Lalonde et al. [1] introduced the concept of utilizing ``weak cues'' extracted from various image segments---the sky region, vertical surfaces, and the ground plane---to estimate sun position and visibility. Their work demonstrated that while any single visual feature in isolation may be uninformative, the probabilistic integration of multiple complementary cues leads to a robust estimate of the sky dome's appearance. The present project adopts this multi-cue philosophy by analyzing both sky gradients and color temperatures to identify solar characteristics, extending the concept to a learned feature-fusion framework.

The ``Photo Sundial'' system proposed by Wu et al. [3] explicitly addressed the issue of unreliable EXIF timestamps in consumer photography. The authors developed a framework that utilizes astronomical theory, linking the sun's apparent position and the camera's viewing direction to calculate the exact time of capture. This study is particularly relevant to the forensic image analysis motivation of the present work and the correction of mislabeled gallery metadata.

Recent advancements by Hold-Geoffroy et al. [2] shifted the paradigm from analytical parametric models to learned high dynamic range (HDR) sky representations. They highlighted the limitations of traditional parametric sky models (e.g., the Hosek--Wilkie model) in representing non-clear-sky conditions such as overcast or partially cloudy weather. By training a convolutional neural network (CNN) to learn sky representations directly from data, they achieved superior results under diverse meteorological conditions.

Unlike previous studies that often treat the date of capture as a static categorical label or ignore it entirely, this project integrates the calendar date directly into the model's feature vector through cyclic encoding. This allows the system to distinguish between identical lighting conditions that occur at different solar angles during the winter and summer solstices. The purpose of this study is to develop a hybrid deep learning framework that transcends simple luminance metrics, combining visual feature extraction with temporal metadata integration for robust, continuous time-of-day regression from consumer sky photographs.


% ========================================================================
%  2. MATERIALS AND METHODS
% ========================================================================
\section{Materials and Methods}

This section describes the complete pipeline for time-of-day estimation, encompassing data collection, preprocessing, feature extraction, model architecture, training strategy, and evaluation protocol. Two parallel implementations were developed: a primary Python/PyTorch pipeline centered on a ConvNeXt backbone, and a secondary MATLAB pipeline utilizing SqueezeNet with classical machine learning baselines.

\subsection{Dataset Construction and Preprocessing}

The dataset comprises outdoor sky photographs collected from the personal galleries of the research team members using consumer-grade smartphone cameras. Each image is required to contain visible sky content and to carry valid EXIF metadata including the \texttt{DateTimeOriginal} tag. Supported image formats include JPEG, PNG, HEIC, and DNG (raw).

For each image, the following metadata fields are extracted from the EXIF header using a dual-library strategy (\texttt{exifread} as the primary parser, with Pillow as a fallback): (i) capture time, parsed from \texttt{DateTimeOriginal}, \texttt{DateTime}, or \texttt{DateTimeDigitized} tags and converted to minutes since midnight ($t \in [0, 1440)$); (ii) calendar date (month, day, year) used to compute the day of year ($d \in [1, 366]$); and (iii) GPS coordinates (latitude and longitude), normalized to $[-1, 1]$. Images lacking valid temporal EXIF data are excluded.

All images are preprocessed through a resolution-standardization pipeline. Original images are downscaled to a maximum dimension of 512 pixels using Lanczos resampling while preserving the aspect ratio. Downscaled images are saved as JPEG files at 90\% quality with EXIF metadata preserved using the \texttt{piexif} library. At training time, images are further processed through a letterbox resize to a fixed $512 \times 512$ canvas with black padding and normalized using ImageNet statistics ($\mu = [0.485, 0.456, 0.406]$, $\sigma = [0.229, 0.224, 0.225]$). For the MATLAB pipeline, images are resized to $224 \times 224$ pixels.

\subsection{Feature Engineering}

\subsubsection{Handcrafted Photometric Features}

A comprehensive set of 77 handcrafted image features is extracted from each photograph. These features are organized as follows: RGB channel statistics (mean and standard deviation per channel, 6 dims), HSV channel statistics (6 dims), CIELAB channel statistics normalized (6 dims), RGB histograms with 8 bins per channel (24 dims), HSV histograms with 8 bins per channel (24 dims), sky-region luminance computed as mean and standard deviation of the value channel in the top third of the frame (2 dims), vertical luminance gradient (1 dim), color temperature proxy as the R/B channel ratio (1 dim), global luminance with mean, standard deviation, and entropy (3 dims), saturation statistics (2 dims), Canny edge density (1 dim), and Laplacian variance as a cloud texture measure (1 dim).

Color space conversions are performed using the \texttt{scikit-image} library. The sky-region luminance features specifically target the upper third of the image frame, which typically contains the sky in landscape photographs and provides the strongest cues about solar position.

\subsubsection{Calendar Metadata Encoding}

Temporal metadata is encoded using cyclic trigonometric transformations to preserve circular continuity:

\begin{equation}
\mathbf{m}_{\text{cal}} = \begin{bmatrix}
    \sin\!\left(\frac{2\pi \cdot \text{month}}{12}\right),
    \cos\!\left(\frac{2\pi \cdot \text{month}}{12}\right),
    \sin\!\left(\frac{2\pi \cdot \text{doy}}{365.25}\right),
    \cos\!\left(\frac{2\pi \cdot \text{doy}}{365.25}\right),
    \frac{\text{lat}}{90},
    \frac{\text{lon}}{180}
\end{bmatrix}^T
\end{equation}

This encoding ensures that December and January are treated as adjacent months. The complete metadata vector has dimensionality $d_{\text{meta}} = 6 + 77 = 83$ when image features are enabled.

\subsection{Target Encoding}

The regression target is encoded as a unit-circle representation to respect the circular topology of the 24-hour clock:

\begin{equation}
\mathbf{y} = \left[\sin\!\left(\frac{2\pi t}{1440}\right),\;
\cos\!\left(\frac{2\pi t}{1440}\right)\right]^T
\end{equation}

where $t$ is the time in minutes since midnight. Decoding is performed via $\hat{t} = \frac{\text{atan2}(\hat{y}_{\sin}, \hat{y}_{\cos}) \bmod 2\pi}{2\pi} \times 1440$. This encoding eliminates the artificial discontinuity at midnight that would cause large gradients in a naive linear regression formulation.

\subsection{Model Architecture}

\subsubsection{Primary Model: ConvNeXt with Metadata Fusion}

The primary model consists of two components: a convolutional feature encoder and a metadata fusion multilayer perceptron (MLP).

A ConvNeXt-Tiny backbone [4], pretrained on ImageNet-1K, serves as the visual feature extractor, producing a 768-dimensional feature vector after global average pooling. All layers up to and including \texttt{features.4} are frozen, while the remaining deeper stages are trainable.

The 768-dimensional image feature vector is concatenated with the 83-dimensional metadata vector to form an 851-dimensional input to a three-layer MLP:

\begin{equation}
\begin{aligned}
\mathbf{h}_1 &= \text{Drop}\!\left(\text{GELU}\!\left(
    \text{LN}\!\left(\mathbf{W}_1 [\mathbf{f}_{\text{img}};
    \mathbf{m}] + \mathbf{b}_1\right)\right)\right) \\
\mathbf{h}_2 &= \text{Drop}\!\left(\text{GELU}\!\left(
    \text{LN}\!\left(\mathbf{W}_2 \mathbf{h}_1 +
    \mathbf{b}_2\right)\right)\right) \\
\hat{\mathbf{y}} &= \mathbf{W}_3 \mathbf{h}_2 + \mathbf{b}_3
\end{aligned}
\end{equation}

where $\mathbf{h}_1 \in \mathbb{R}^{384}$, $\mathbf{h}_2 \in \mathbb{R}^{192}$, $\hat{\mathbf{y}} \in \mathbb{R}^{2}$, LN denotes LayerNorm, and the dropout rate is $p = 0.1555$.

\subsubsection{MATLAB Pipeline: SqueezeNet with Baselines}

A parallel MATLAB implementation employs a multi-input CNN based on SqueezeNet. The image branch uses SqueezeNet with fire modules 1--6 frozen and 7--9 fine-tuned, producing a 512-dimensional feature vector. The date branch maps 4-dimensional cyclic date features through FC(32)$\rightarrow$BN$\rightarrow$ReLU$\rightarrow$FC(64)$\rightarrow$ReLU. The concatenated 576-dimensional vector passes through FC(256)$\rightarrow$BN$\rightarrow$ReLU$\rightarrow$Dropout(0.4)$\rightarrow$FC(64)$\rightarrow$ReLU$\rightarrow$FC(1) for regression.

Two classical baselines are trained on a 149-dimensional feature vector (145-dim classical features + 4-dim cyclic date): a Random Forest with 100 trees (\texttt{TreeBagger}) and a Support Vector Regressor with RBF kernel (\texttt{fitrsvm}).

\subsection{Training Strategy}

\subsubsection{Loss Function and Evaluation Metric}

The model is trained using mean squared error loss on the sine--cosine encoded targets:

\begin{equation}
\mathcal{L} = \frac{1}{N} \sum_{i=1}^{N}
    \left\| \hat{\mathbf{y}}_i - \mathbf{y}_i \right\|_2^2
\end{equation}

The primary evaluation metric is the cyclic mean absolute error (MAE) in minutes:

\begin{equation}
\text{MAE}_{\text{cyc}} = \frac{1}{N} \sum_{i=1}^{N}
    \min\!\left(|\hat{t}_i - t_i|,\; 1440 - |\hat{t}_i - t_i|\right)
\end{equation}

\subsubsection{Data Augmentation}

A heavy augmentation regime is employed: random resized crop (scale $[0.8, 1.0]$), random horizontal flip ($p = 0.5$), RandAugment (3 operations, magnitude 12), random erasing ($p = 0.25$, area scale $[0.02, 0.15]$), Mixup ($\alpha = 0.1776$), and Gaussian label noise ($\sigma = 0.0437$).

\subsubsection{Optimization}

The optimizer is 8-bit AdamW [5] with hyperparameters determined via Optuna [6] Bayesian optimization using Tree-structured Parzen Estimator sampling: learning rate $\eta = 4.42 \times 10^{-4}$, weight decay $\lambda = 0.0373$, cosine annealing schedule with $\eta_{\min} = 1.83 \times 10^{-5}$, 80 training epochs, batch size of 8, and gradient clipping at max norm 1.0. Mixed-precision training (FP16), \texttt{torch.compile()} graph optimization, and NHWC memory format are enabled for throughput gains.

The hyperparameter search space includes learning rate ($[1.5\text{e-}4, 4.5\text{e-}4]$, log), weight decay ($[1.5\text{e-}2, 6.5\text{e-}2]$, log), cosine $\eta_{\min}$ ($[5.0\text{e-}6, 2.0\text{e-}5]$, log), dropout ($[0.10, 0.25]$), label noise ($[0.02, 0.05]$), Mixup $\alpha$ ($[0.10, 0.30]$), augmentation magnitude (\{moderate, heavy\}), and freeze depth (\{features.4, features.6\}). The search uses a median pruner with 15 total trials.

\subsubsection{Class Imbalance Mitigation}

A weighted random sampler addresses the non-uniform distribution of capture times. Per-sample weights are computed as the inverse frequency of the corresponding hourly bin, with extreme weights capped at $10\times$ the median weight.

\subsection{Cross-Validation Protocol}

Both pipelines employ 5-fold cross-validation using stratified random shuffling with a fixed seed of 42. In each fold, 80\% of the data is used for training and 20\% for validation. In the MATLAB pipeline, the CNN training set is further split with a 15\% holdout for early stopping (patience of 7), ensuring the test fold is never used for model selection.

\subsection{Ensemble Inference}

The ensemble module aggregates predictions from all five fold-specific checkpoints by averaging raw sine--cosine outputs and re-normalizing onto the unit circle:

\begin{equation}
\hat{\mathbf{y}}_{\text{ens}} = \frac{1}{K} \sum_{k=1}^{K}
    \hat{\mathbf{y}}_k, \qquad
\tilde{\mathbf{y}} = \frac{\hat{\mathbf{y}}_{\text{ens}}}
    {\|\hat{\mathbf{y}}_{\text{ens}}\|_2}
\end{equation}

where $K = 5$. This circular averaging avoids artifacts from naively averaging decoded minute values near midnight.


% ========================================================================
%  3. CONCLUSIONS
% ========================================================================
\section{Conclusions}

This paper presented a comprehensive framework for estimating the time of day from single sky photographs by fusing deep convolutional features with handcrafted photometric descriptors and cyclically encoded calendar metadata. The proposed methodology addresses the fundamental limitations of brightness-based heuristics used in conventional image organization systems by incorporating multi-modal cues---color temperature, sky gradients, edge density, seasonal solar positioning, and geographic coordinates---into a unified regression pipeline.

The cyclic sine--cosine encoding scheme for both the target variable and temporal metadata features proved essential for modeling the circular nature of the 24-hour clock and the annual calendar cycle, eliminating artificial discontinuities that would otherwise degrade regression performance. The integration of 77-dimensional handcrafted features alongside the ConvNeXt-Tiny backbone's learned representations provided complementary information channels that enhanced estimation accuracy beyond what either modality could achieve independently.

The dual-pipeline implementation---a primary Python/PyTorch system and a parallel MATLAB-based system with classical baselines---enabled a meaningful comparative analysis between modern deep learning approaches and established machine learning techniques for this regression task. The Optuna-driven hyperparameter optimization ensured that the reported configurations represent near-optimal operating points for the given dataset and architecture.

Future work may extend this framework in several directions: incorporating attention mechanisms to automatically identify informative image regions, leveraging sun position estimation from shadow analysis as an auxiliary prediction target, expanding the dataset to include diverse geographic locations and climatic conditions to improve generalization, and deploying the trained model as a lightweight mobile application for real-time image metadata correction.


% ========================================================================
%  AUTHOR CONTRIBUTION STATEMENT
% ========================================================================
\section*{Author Contribution Statement}

The contributions of each author to this study are as follows. \textbf{Author 1} (A.G. Efe, 33\%): literature review, MATLAB pipeline development including SqueezeNet CNN and classical baseline implementations, data collection. \textbf{Author 2} (D. Parlaky\i ld\i z, 33\%): Python/PyTorch pipeline development, ConvNeXt model architecture design, hyperparameter optimization with Optuna, writing. \textbf{Author 3} (S. S\"unb\"ul, 34\%): dataset construction and preprocessing pipeline, feature engineering, cross-validation framework, ensemble inference module, critical review.


% ========================================================================
%  REFERENCES
% ========================================================================
\section*{References}
\begin{enumerate}[label={[\arabic*]}, leftmargin=*, itemsep=2pt, parsep=0pt, font=\small]

\item Lalonde, J.-F., Efros, A.A., Narasimhan, S.G. 2012. Estimating the Natural Illumination Conditions from a Single Outdoor Image, \textit{International Journal of Computer Vision}, Vol.~98, No.~2, pp.~123--145. DOI: 10.1007/s11263-011-0501-8

\item Hold-Geoffroy, Y., Sunkavalli, K., Hadap, S., Gambaretto, E., Lalonde, J.-F. 2019. Deep Outdoor Illumination Estimation. \textit{Proceedings of the IEEE/CVF Conference on Computer Vision and Pattern Recognition}, pp.~7312--7321. DOI: 10.1109/CVPR.2019.00749

\item Wu, L., Ganesh, A., Shi, B., Matsushita, Y., Wang, Y., Ma, Y. 2011. Robust Photometric Stereo via Low-Rank Matrix Completion and Recovery. \textit{Proceedings of the Asian Conference on Computer Vision}, Springer.

\item Liu, Z., Mao, H., Wu, C.-Y., Feichtenhofer, C., Darrell, T., Xie, S. 2022. A ConvNet for the 2020s. \textit{Proceedings of the IEEE/CVF Conference on Computer Vision and Pattern Recognition}, pp.~11976--11986. DOI: 10.1109/CVPR52688.2022.01167

\item Dettmers, T., Lewis, M., Belkada, Y., Zettlemoyer, L. 2022. 8-bit Optimizers via Block-wise Quantization. \textit{Proceedings of the International Conference on Learning Representations}.

\item Akiba, T., Sano, S., Yanase, T., Ohta, T., Koyama, M. 2019. Optuna: A Next-generation Hyperparameter Optimization Framework. \textit{Proceedings of the ACM SIGKDD International Conference on Knowledge Discovery and Data Mining}, pp.~2623--2631. DOI: 10.1145/3292500.3330701

\end{enumerate}

\end{document}
